% VI25_FULL_SOLUTION_REFINEMENT_V1
\subsection*{VI.25.P2: Affine fiber formula}
\begin{problem}\label{prob:vi25-p02}
Let \(A\to B\) and \(\mathfrak p\in\Spec A\).  Prove
\[
(\Spec B)_{\mathfrak p}\cong\Spec(B\otimes_A\kappa(\mathfrak p)).
\]
\end{problem}

\ifdefined\FullProblemDossiers
\paragraph{Why this problem matters.}
This formula is the computational engine of the chapter: every affine fiber calculation reduces to one tensor product.

\paragraph{Useful ideas before starting.}
Apply the affine fiber-product theorem to \(\Spec B\to\Spec A\leftarrow\Spec\kappa(\mathfrak p)\), then rewrite the tensor product as a localization of \(B/\mathfrak pB\).

\paragraph{Guided hints.}
\begin{enumerate}
\item Use \(X_{\mathfrak p}=X\times_{\Spec A}\Spec\kappa(\mathfrak p)\).
\item Recall that \(\kappa(\mathfrak p)\cong (A/\mathfrak p)_{A\setminus\mathfrak p}\).
\item Explain the two algebraic steps: quotient by \(\mathfrak p\), then invert \(A\setminus\mathfrak p\).
\end{enumerate}

\begin{solution}
By definition,
\[
(\Spec B)_{\mathfrak p}
=
\Spec B\times_{\Spec A}\Spec\kappa(\mathfrak p).
\]
All three schemes are affine, so VI/23 gives
\[
\Spec B\times_{\Spec A}\Spec\kappa(\mathfrak p)
\cong
\Spec\!\bigl(B\otimes_A\kappa(\mathfrak p)\bigr).
\]
This proves the stated formula.

To see what the ring means, put \(S=A\setminus\mathfrak p\).  Since
\[
\kappa(\mathfrak p)
\cong
S^{-1}(A/\mathfrak p),
\]
associativity of tensor products and compatibility with localization give
\[
B\otimes_A\kappa(\mathfrak p)
\cong
S^{-1}(B/\mathfrak p B)
\cong
S^{-1}B/\mathfrak p S^{-1}B.
\]
Thus the fiber first forces every element of \(\mathfrak p\) to vanish and then makes every element of \(A\setminus\mathfrak p\) invertible.  This is why a fiber is a specialization of equations rather than merely a set-theoretic inverse image.
\end{solution}

\paragraph{What happened mathematically.}
Scheme-theoretic specialization is encoded by the tensor product with the residue field.

\paragraph{Extensions.}
Every later affine fiber computation in the chapter is an instance of this formula.

\paragraph{Provenance.}
Canonical affine-fiber theorem derived from VI/23 and VI/24.
\else
\paragraph{Provenance.} Canonical affine-fiber theorem derived from VI/23 and VI/24.
\fi
